\section{Per-gate exact and approximate counting}
\label{sec:unified}

Let $C$ be a single-output Boolean DAG. Its gates have ordered input
lists, possibly with repetitions. Put $S=n+w+2$ and
$b=\lceil S^{9/10}\rceil$. A gate of fan-in $r_v>b$ is large; all
other gates are small. The allowed types are
\[
\begin{aligned}
\THR &: \ind{\sum_jw_jy_j\ge\theta},\\
\ETHR &: \ind{\sum_jw_jy_j=\theta},\\
\MOD_m &: \ind{\sum_jw_jy_j\equiv\theta\pmod m},
\end{aligned}
\]
where all parameters are integers and $m\ge1$. Let $L$ be the length
of the complete input description, including the mode assignment and
the certificates below, and fix $a\ge1$. A small gate can instead
be an arbitrary Boolean function supplied by an evaluator. When any
gate is hashed, require $L\le S^a$ and evaluator time at most $S^a$.
For an exact-only assignment, $L$ may be arbitrary and evaluator time
at most $(S+L)^a$ is sufficient.

For a large threshold or exact-threshold gate, let $B_v$ be the maximum
bit length of $|w_j|$, at least one. For a modular gate let
$B_v=\max\{1,\bits(m-1)\}$. Each large gate receives one of the
following modes. In the factorization mode, preprocessing verifies the
dimensions and entry bounds of the certificate and checks each row of
$w=C_vg_v$ by integer arithmetic. A malformed certificate is rejected;
the verification costs polynomial time in $S+L$.
\begin{center}\small
\renewcommand{\arraystretch}{1.15}
\begin{tabular}{@{}P{.16\linewidth}P{.22\linewidth}P{.37\linewidth}P{.15\linewidth}@{}}
\toprule
Mode & Gate types & Cost $\sigma_v$ & Error\\
\midrule
Exact slices & All three & $B_v$ & $0$\\
Exact factorization & All three & $\kappa_v=d_v(\beta_v+\lceil\log_2(r_v+1)\rceil+2)$ & $0$\\
Hashed comparison & $\THR$ & $h_v\lceil\log_2S\rceil$, $h_v\ge1$ & $2^{-h_v}$\\
Hashed equality & $\ETHR$, $\MOD_m$ & $h_v$, $h_v\ge1$ & $2^{-h_v}$\\
\bottomrule
\end{tabular}
\end{center}
For factorization, $C_v$ is an $r_v\times d_v$ integer matrix with
entries of absolute value less than $2^{\beta_v}$, $\beta_v\ge1$,
and $g_v\in\Z^{d_v}$. Its coordinates can be large; their full bit
length is charged to $L$.

\begin{theorem}[Per-gate counting, conditional on \HTC]
\label{thm:unified}
Assume \HTC. Fix $a,c\ge1$. Suppose the assigned modes satisfy
\begin{equation}\label{eq:budget}
 \sum_{v\text{ large}}r_v\left\lceil\frac{\sigma_v}{b}\right\rceil
 \le cS,
 \qquad \sigma_v\le cb
 \quad\text{for every factorized or hashed large gate}.
\end{equation}
Let $\mathcal H$ be the set of hashed gates. If
$\mathcal H\ne\varnothing$, under $L\le S^a$, in
$\poly_{a,c}(S)2^{O_{a,c}(g(S))}$ time and polynomial space a
deterministic algorithm outputs a rational $\widehat p$ with
\[
 \left|\widehat p-\Pr_x[C(x)=1]\right|
 \le\sum_{v\in\mathcal H}2^{-h_v}.
\]
If $\mathcal H=\varnothing$, the restriction $L\le S^a$ is unnecessary:
the algorithm outputs the exact accepting-input count in
$\poly_{a,c}(S+L)2^{O_{a,c}(g(S))}$ time and
$\poly_{a,c}(S+L)$ space, under the exact-only evaluator condition above.
\end{theorem}

\subsection{Exact sliced gadgets}

A small gate has the single factor $[z=f(y)]$ of degree $r_v+1$.
Now fix a large gate with fan-in $r$, width $B$, and let
$K=\lceil B/b\rceil$. For $\THR$ or $\ETHR$, replace signed weights
by $u_j=|w_j|$ and the corresponding literal $\xi_j$ as above. Write
$T=\theta+\sum_{w_j<0}|w_j|$ and
$\Sigma=\sum_ju_j\xi_j$. We have
$0\le\Sigma<r2^B$. For a threshold gate clamp $T<0$ to zero; for an
exact-threshold gate replace $T<0$ by $r2^B$. In either case clamp
$T>r2^B$ to $r2^B$. These operations preserve the Boolean function.
For a modular gate, replace $w_j$ by $u_j=w_j\bmod m$ and $\theta$
by $\theta'=\theta\bmod m$, with $\xi_j=y_j$. Then $u_j<2^B$ and
$0\le\Sigma<rm$.

Decompose each $u_j$ into base-$2^b$ digits $u_j^{(k)}$, $0\le k<K$.
For each digit position compute
\[
 Y^{(k)}=\sum_ju_j^{(k)}\xi_j,\qquad
 \Sigma=\sum_{k=0}^{K-1}2^{kb}Y^{(k)}.
\]
Partition the $r$ inputs into $t=\lceil r/b\rceil\le2r/b$ blocks.
A chain for each $k$ reads one block at a time and checks the exact
partial-sum update, passing a state of
$\mu=b+\lceil\log_2(r+1)\rceil$ bits. Each input is read $K$ times.
The chain has $O(r)$ wires per slice, degree $b+2\mu$, and unique
states on every original input.

The final arithmetic relates the slice sums by short carries. For
nonnegative constants $A,T'$ occurring below put
\[
 K^*=\left\lceil\frac{B+\lceil\log_2(r+2)\rceil}{b}\right\rceil.
\]
They are less than $2^{K^*b}$. Let $A_k,T'_k$ be their base-$2^b$
digits and $Y^{(k)}=0$ for $k\ge K$. Starting with
$c_{-1}=0,e_{-1}=1$, enforce
\begin{align}
 z_k&=Y^{(k)}+A_k-T'_k+c_{k-1},\nonumber\\
 c_k&=\left\lfloor z_k/2^b\right\rfloor,
 &e_k&=e_{k-1}\wedge[z_k\equiv0\pmod{2^b}].
 \label{eq:carry}
\end{align}
Induction shows
\[
 \Sigma+A-T'=c_{K^*-1}2^{K^*b}+R_0,
 \qquad 0\le R_0<2^{K^*b}.
\]
Thus the difference is nonnegative exactly when the last carry is
nonnegative, and it is zero exactly when that carry is zero and the
last zero flag is one. Every carry lies in $[-1,r+1]$: if the previous
one does, then $-2^b\le z_k\le(r+1)2^b$. An offset encoding therefore
uses $\lceil\log_2(r+3)\rceil$ bits. With $A=0,T'=T$ the last carry
and flag give the required $\THR$ or $\ETHR$ output.

For $\MOD_m$, put $A=m-\theta'$ and introduce a quotient code
$q\in\{0,\ldots,r\}$. Compute two instances of~\eqref{eq:carry},
with $T'=qm$ and $T'=(q+1)m$. The terminal factor requires
\[
 0\le\Sigma+(m-\theta')-qm<m
\]
and outputs one exactly when its left side is zero. There is precisely
one legal $q$, namely $\lfloor(\Sigma+m-\theta')/m\rfloor$; invalid
binary codes are rejected. Every carry factor can read the short
$q$ code and compute its needed digits locally. The same quotient
gadget includes $m=1$ without a separate case.

The slice chains, carries, quotient copies, and output use $O(rK)$
wires, degree $O(b+\log S)$, and polynomial-time local entry programs.
The bound $K^*\le K+1$ holds for sufficiently large $S$; more generally
$K^*=K+O(1)$ suffices. The $O(K\log S)$ carry and quotient wires
are absorbed by $O(Kr)$ because $r>b$. All internal states and the
output are unique once the input is fixed. This establishes the exact
sliced mode.

\subsection{Exact factorization gadgets}

For $w=Cg$, write $\kappa=\kappa_v$ for this gate's state width.
The information to transmit across blocks is a vector,
rather than the weighted sum itself. After block $i$ store
\[
 Y_i=\sum_{j\text{ in blocks }1,\ldots,i}y_jC_{j,\cdot}\in\Z^d.
\]
Each coordinate has absolute value less than $r2^\beta$. A fixed
signed code of $\beta+\lceil\log_2(r+1)\rceil+2$ bits suffices, so
the complete vector uses $\kappa$ bits. Block factors enforce the
updates. The last factor checks the desired comparison, equality,
or congruence of $g\cdot Y_t$ with $\theta$. The supplied identity
$w=Cg$ proves correctness. Updates and output are unique. Before
equality expansion and subdivision, the logical wire count satisfies
\begin{align*}
 r+t\kappa+1
 &\le r+(r/b+1)\kappa+1\\
 &\le r+cr+cb+1\le(2c+2)r,
\end{align*}
using $t=\lceil r/b\rceil\le r/b+1$, $\kappa\le cb$, and $r>b$.
Thus there are $O_c(r)$ wires for every fixed $c\ge1$, and degree at most
$b+2\kappa+O(1)=O_c(b)$ when $\kappa\le cb$.

For example, if a gate's weights have at most $d$ distinct values,
one may take $C$ to be their indicator matrix and $g$ the list of
values. This gives $\beta=1$ and $\kappa=O(d\log S)$ independently
of the bit lengths of those values. The factorization is part of the
input; no search for an optimal one is used.

\subsection{Hashed equality and congruence}

For an exact-threshold gate set
$\Sigma=\sum_jw_jy_j$ and $\mathcal T=\{\theta\}$. For a modular
gate use reduced weights and the target set
\[
 \mathcal T=\{\theta'+qm:0\le q<r\}.
\]
Because $\Sigma\in[0,rm)$, membership in this set is exactly the
original modular predicate. Given a prime $p>S$, replace either
predicate by
\[
 f_p(y)=\ind{\Sigma\bmod p\in\mathcal T\bmod p}.
\]
Every true input remains true. For a false input, a false positive
requires a prime divisor of some nonzero $\Sigma-\tau$,
$\tau\in\mathcal T$. Each such integer has bit length at most
$S^{\lceil a\rceil}+2\lceil\log_2S\rceil+3$. Put
\[
 N_v=r(S^{\lceil a\rceil}+2\lceil\log_2S\rceil+3)
\]
and let $K_v$ be the least power of two at least $N_v2^{h_v}$.
Selecting among the $K_v$ smallest primes greater than $S$ gives error
at most $2^{-h_v}$. Their common bit bound is
$\pi_v=h_v+O_a(\log S)$.

A block chain passes the partial sum modulo $p$ in $\pi_v$ bits.
Its last factor tests membership in the at-most-$r$ reduced targets.
Reject residue codes outside $[0,p)$. For each fixed prime this is a
unique-extension circuit relation, with $O_{a,c}(r)$ wires and degree
$O_{a,c}(b)$ if $h_v\le cb$. In particular the end-state contribution
is bounded by $O_{a,c}(r)$; it need not be at most $r$ when $c>1$.

\subsection{Hashed comparisons}

For a large $\THR$ gate, use the compact comparison realization
\eqref{eq:compact-edges}--\eqref{eq:compact-degree}, grouping $b$
input occurrences per party and choosing error $2^{-h_v}$. Its
parameters satisfy
\[
 J=O_a(\log S),\quad \pi_v=O_a(h_v+\log S),\quad
 \lambda=O(\log S),\quad t\le2r/b.
\]
If $h_v\lceil\log_2S\rceil\le cb$, then
$J\pi_v=O_{a,c}(b)$ and $J^2\lambda=o(b)$. The edge count is
$O_{a,c}(r)$ and the degree is $O_{a,c}(b)$. For the terminal
$J\pi_v$ term, use $r>b$. The seed length is
$O_a(h_v\log S+\log^2S)=O_{a,c}(b)$.
The comparison's pointwise error and the functional property are the
ones proved in Section~\ref{sec:protocol}.

\subsection{Assembling the network}

\begin{proof}[Proof of Theorem~\ref{thm:unified}]
For every hashed gate use the required prefix of a common seed of
length $R$, where $R$ is the maximum of the gate seed lengths and is
zero if no gate is hashed. Both kinds of hashing above give
$R=O_{a,c}(S^{9/10})$. A uniform seed has uniform prefixes, so each
gate retains its own error guarantee. Generate the selected primes
by index, as described above. Per seed the extra time is
$2^{O_{a,c}(S^{9/10})}$ and the space is polynomial.

Replace each gate by its gadget. Connect every occurrence of each
original signal by the equality construction of
Section~\ref{sec:interface}, including a factor two for each unused
input, and pin the final output. Internal signals read by several
factors, including the MOD quotient code and the comparison protocol's
branch history, also use equality trees; these copies are already
included in the gadget port counts. For a fixed seed the resulting
network counts precisely the accepting assignments of $C_\rho$.
Its edges, including equality and bipartite subdivision, are at most
\[
 O_{a,c}\left(n+w+
 \sum_{v\text{ large}}r_v\lceil\sigma_v/b\rceil\right)=O_{a,c}(S).
\]
Its degree is at most $D_{a,c}S^{9/10}$ for a fixed $D_{a,c}$ and its
vertex count is $O_{a,c}(S)$. Its description length and entry costs
are polynomial in $S+L$, hence polynomial in $S$ in the hashed case.
Use~\eqref{eq:rescale} with $s=\lceil K_{a,c}S\rceil$,
$P_0=n+1$, and $\ell=1$. All the hypotheses of \HTC\ hold.

Recover the count $Z_\rho\in[0,2^n]$ exactly. If $R=0$, this is the
answer. Otherwise form~\eqref{eq:capp-output}. Applying
\eqref{eq:gate-error} and averaging over $x$ gives the stated error.
The number of seeds and the prime-generation cost together add
$2^{O_{a,c}(S^{9/10})}$, absorbed by $2^{O_{a,c}(g(S))}$.
All calls are serial, and the accumulator uses $n+R+1$ bits.

In the exact-only case there are no primes to generate and $R=0$.
All arithmetic on weights, thresholds, moduli, and supplied certificates
uses polynomial time and space in $S+L$. The edge, degree, vertex,
and precision bounds above depend on the mode budget, not on an upper
bound for $L$. Retain the original input in the tensor description,
so its length $\mathcal L$ satisfies
$L\le\mathcal L\le\poly_{a,c}(S+L)$. Thus the local evaluation
bound is also polynomial in $S+\mathcal L$, and the same \HTC\ call
gives the stated arbitrary-description bound.
\end{proof}

\begin{proof}[Proof of Theorem~\ref{thm:capp}, including high precision]
Use Theorem~\ref{thm:unified} with the fixed description exponent
$\lceil a\rceil+2$ to include generated mode parameters; all resulting
constants still depend only on $a$.
Assign hashed-comparison mode to every large gate and choose
\[
 h_v=\left\lceil\log_2(2S/\eps)\right\rceil.
\]
Then $\sum_v2^{-h_v}\le\eps/2$, and at the lower endpoint
$h_v\lceil\log_2S\rceil\le2b$ for all sufficiently large $S$.
Each summand in~\eqref{eq:budget} is at most $2r_v$; taking a fixed
$c\ge2$ proves the required budget. For a parameter-free endpoint
implementation one can use the integer target
\[
 h_v=\left\lceil\frac{\lceil S^{9/10}\rceil}
 {\lfloor\log_2(S+16)\rfloor}\right\rceil
 +\lceil\log_2(2S)\rceil+2,
\]
which gives at least the requested accuracy and still obeys the same
constant budget. Finite small $S$ are enumerated. Apply
Theorem~\ref{thm:unified}.
\end{proof}

\begin{corollary}[Exact specializations, conditional on \HTC]
\label{cor:exact}
Assume \HTC\ and the exact-only evaluator condition above; let $L$ be
the arbitrary complete description length.
If all large gates are assigned exact slices and
\[
 \sum_{v\text{ large}}r_vB_v\le c_1S^{19/10}
\]
for a fixed $c_1$, exact counting takes
$\poly_{a,c_1}(S+L)2^{O_{a,c_1}(g(S))}$ time and
$\poly_{a,c_1}(S+L)$ space.
The same holds when every large gate has a supplied factorization
with $\kappa_v\le c_1b$. Small gates may have arbitrary parameters
with their full lengths charged to $L$. In particular, this recovers exact
counting for bounded-fan-in circuits at the allowed degree scale and
for large gates whose weight bit lengths are $O(S^{9/10})$; the
modular sliced mode also allows modulus bit lengths $O(S^{9/10})$.
\end{corollary}
\begin{proof}
For slicing use
$\lceil B_v/b\rceil\le1+B_v/b$ and $\sum_vr_v\le w<S$.
For factorization use
$\lceil\kappa_v/b\rceil\le\lceil c_1\rceil$.
Theorem~\ref{thm:unified} has no hashed gate in either case.
\end{proof}

\subsection{Relations between the specializations}

For integer weights, an equality gate has the identity
\begin{equation}\label{eq:equality-thresholds}
 \ind{\sum_jw_jy_j=\theta}
 =\ind{\sum_jw_jy_j\ge\theta}
 \wedge\neg\ind{\sum_jw_jy_j\ge\theta+1}.
\end{equation}
This uses the same weight vector in two threshold gates and only a
constant-factor increase in wires. Ordinary unweighted MOD gates of
fixed modulus likewise have $O(r)$-wire Boolean, and hence threshold,
simulations: pass the running residue in a constant number of bits.
They therefore fit the threshold-circuit lower-bound framework of
Section~\ref{sec:lower}. For a fixed modulus, arbitrary weights can
also be reduced modulo that modulus before this simulation. The
modular cases requiring the bit-length or certificate accounting here
are predicates with growing modulus and potentially large reduced
weights.

The sliced budget is an aggregate condition even for pure THR
circuits, whereas the factorization condition bounds each supplied
state width $\kappa_v$. For example, a large gate with
$r_v=\Theta(S^{.95})$ and $B_v=\Theta(S^{.94})$ contributes
$\Theta(S^{1.89})$ to the sliced sufficient condition, while its
one-dimensional certificate $C=w$, $g=1$ has width
$\kappa_v=B_v+O(\log S)$, above the factorization cap.
A shorter certificate is a separate property of its weights. Thus
the two sufficient conditions are not identified, including on pure
THR circuits. With $L\le S^a$ they are specializations of the
polynomial-description statement; their arbitrary-$L$ versions use
its exact-only clause.
